% =====================================================================
% ACM acmart sigconf 会议论文最小可编译骨架
% 编译: latexmk -pdf acm_sigconf.tex            (pdfLaTeX)
% 双盲审稿: \documentclass[sigconf,review,anonymous]{acmart}
% 最终版须填 venue 提供的 rights/DOI/ISBN, 否则出红色版权提示框
% CCS concepts 从 https://dl.acm.org/ccs 生成后粘贴
% =====================================================================
\documentclass[sigconf]{acmart}

% --- 版权与会议信息 (占位, 投稿时换成 venue 提供值) ----------------
\setcopyright{acmlicensed}
\copyrightyear{2026}
\acmYear{2026}
\acmDOI{XXXXXXX.XXXXXXX}
\acmConference[Conf '26]{The ACM Conference}{June 03--05, 2026}{City, Country}
\acmISBN{978-1-4503-XXXX-X/26/06}

\begin{document}

\title{Paper Title Here}
\subtitle{Optional Subtitle}

\author{First Author}
\affiliation{%
  \institution{Institution Name}
  \city{City}
  \country{Country}}
\email{first@example.com}

\author{Second Author}
\affiliation{%
  \institution{Institution Name}
  \city{City}
  \country{Country}}
\email{second@example.com}

% 页眉短作者/短标题
\renewcommand{\shortauthors}{Author et al.}

\begin{abstract}
This is the abstract. State problem, method, results and contributions in one paragraph.
\end{abstract}

% --- CCS 分类 (从 dl.acm.org/ccs 生成完整 CCSXML 后替换) ----------
\begin{CCSXML}
<ccs2012>
<concept>
<concept_id>10010147.10010178</concept_id>
<concept_desc>Computing methodologies~Artificial intelligence</concept_desc>
<concept_significance>500</concept_significance>
</concept>
</ccs2012>
\end{CCSXML}
\ccsdesc[500]{Computing methodologies~Artificial intelligence}

\keywords{keyword1, keyword2, keyword3}

\maketitle

\section{Introduction}
Body text. Citation example~\cite{ref1}. See Fig.~\ref{fig:arch} and Eq.~\eqref{eq:main}.

\section{Method}
\begin{equation}\label{eq:main}
  y = f(x) + \epsilon
\end{equation}

\begin{figure}[t]
  \centering
  % \includegraphics[width=\columnwidth]{fig1.pdf}
  \rule{0.6\columnwidth}{3cm}  % 占位, 有图时替换
  \caption{System architecture.}
  \Description{Architecture diagram.}  % acmart 要求图的可访问性描述
  \label{fig:arch}
\end{figure}

\section{Conclusion}
Summary.

% --- 参考文献 ------------------------------------------------------
% 推荐 bibtex: 把条目放 refs.bib
\bibliographystyle{ACM-Reference-Format}
% \bibliography{refs}
% 演示用内联 (无 refs.bib 时):
\begin{thebibliography}{1}
\bibitem{ref1} A. Author. 2024. Title of paper. In \textit{Proc. of Conf.} 1--10.
\end{thebibliography}

\end{document}
